\chapter{2024年上海卷（春）}

\section{填空题}

\begin{problem}
函数 \(y=\log_2 x\) 的定义域为\fillinblank{} .
\end{problem}

\begin{problem}
直线 \(x-y+1=0\) 的倾斜角为\fillinblank{} .
\end{problem}

\begin{problem}
已知 \(\frac{z}{1+\i}=i\)，则 \(z=\)\fillinblank{} .
\end{problem}

\begin{problem}
\(\left(x-1\right)^6\) 展开式中 \(x^4\) 的系数为\fillinblank{} .
\end{problem}

\begin{problem}
三角形 \(ABC\) 中，\(BC=2\)，\(A=\frac{\pi}{3}\)，\(B=\frac{\pi}{4}\)，则 \(AB=\)\fillinblank{} .
\end{problem}

\begin{problem}
已知 \(ab=1\)，\(4a^2+9b^2\) 的最小值为\fillinblank{} .
\end{problem}

\begin{problem}
数列 \(\{a_n\}\)，\(a_n=n+c\)，\(S_7<0\)，\(c\) 的取值范围为\fillinblank{} .
\end{problem}

\begin{problem}
三角形三边长为 \(5\)，\(6\)，\(7\)，则以边长为 \(6\) 的两个顶点为焦点，过另外一个顶点的双曲线的离心率为\fillinblank{} .
\end{problem}

\begin{problem}
已知 \(f(x)=x^2\)，\(g(x)=\displaystyle\begin{cases}
f(x), & x\ge 0,\\
-f(-x), & x<0,
\end{cases}\) 求 \(g(x)\le 2-x\) 的 \(x\) 的取值范围\fillinblank{} .
\end{problem}

\begin{problem}
在棱柱 \(ABCD-A_1B_1C_1D_1\) 中，底面 \(ABCD\) 为平行四边形，\(\left|AA_1\right|=3\)，\(\left|BD\right|=4\)，\(\overrightarrow{AB_1}\cdot\overrightarrow{BC}-\overrightarrow{AD_1}\cdot\overrightarrow{DC}=5\)，设异面直线 \(AA_1\) 与 \(BD\) 的夹角为 \(\theta\)，则 \(\cos\theta=\)\fillinblank{} .
\end{problem}

\begin{problem}
正方形草地 \(ABCD\) 边长 \(1.2\)，\(E\) 到 \(AB\)，\(AD\) 距离为 \(0.2\)，\(F\) 到 \(BC\)，\(CD\) 距离为 \(0.4\)，有个圆形通道经过 \(E\)，\(F\)，且经过 \(AD\) 上一点，求圆形通道的周长\fillinblank{} .（精确到 \(0.01\)）

\bitmapfigure[max width=0.28\linewidth]{img/2024/shanghai_spring/q11_fig1.png}
\end{problem}

\begin{problem}
\(a_1=2\)，\(a_2=4\)，\(a_3=8\)，\(a_4=16\)，任意 \(b_1\)，\(b_2\)，\(b_3\)，\(b_4\in \R\)，满足 \(\{a_i+a_j\mid 1\le i<j\le 4\}=\{b_i+b_j\mid 1\le i<j\le 4\}\)，求有序数列 \(\{b_1,b_2,b_3,b_4\}\) 有\fillinblank{} 对.
\end{problem}

\section{单选题}

\begin{problem}
已知 \(a\)，\(b\)，\(c\in \R\)，\(b>c\)，下列不等式恒成立的是（\quad）.

\choices
    {\(a+b^2>a+c^2\)}
    {\(a^2-c>a^2-b\)}
    {\(ab^2>ac^2\)}
    {\(a^2b>a^2c\)}
\end{problem}

\begin{problem}
空间中有两个不同的平面 \(\alpha\)，\(\beta\) 和两条不同的直线 \(m\)，\(n\)，则下列说法中正确的是（\quad）.

\choices
    {若 \(\alpha\perp\beta\)，\(m\perp\alpha\)，\(m\perp n\)，则 \(n\perp\beta\)}
    {若 \(\alpha\perp\beta\)，\(m\perp\alpha\)，\(n\perp\beta\)，则 \(m\perp n\)}
    {若 \(\alpha\parallel\beta\)，\(m\parallel\alpha\)，\(n\parallel\beta\)，则 \(m\parallel n\)}
    {若 \(\alpha\parallel\beta\)，\(m\parallel\alpha\)，\(m\parallel n\)，则 \(n\parallel\beta\)}
\end{problem}

\begin{problem}
有四种礼盒，前三种里面分别仅装有中国结，记事本，笔袋，第四个礼盒里面三种礼品都有，现从中任选一个盒子，设事件 \(A\)：所选盒中有中国结，事件 \(B\)：所选盒中有记事本，事件 \(C\)：所选盒中有笔袋，则（\quad）.

\choices
    {事件 \(A\) 与事件 \(B\) 互斥}
    {事件 \(A\) 与事件 \(B\) 相互独立}
    {事件 \(A\) 与事件 \(B\cup C\) 互斥}
    {事件 \(A\) 与事件 \(B\cap C\) 相互独立}
\end{problem}

\begin{problem}
现定义如下：当 \(x\in(n,n+1)\) 时（\(n\in \N\)），若 \(f(x+1)=f'(x)\)，则称 \(f(x)\) 为延展函数. 已知当 \(x\in(0,1)\) 时，\(g(x)=\e^x\)，且 \(h(x)=x^{10}\)，且 \(g(x)\)，\(h(x)\) 均为延展函数，则以下结论（\quad）.

\choices
    {（1）（2）都成立}
    {（1）（2）都不成立}
    {（1）成立（2）不成立}
    {（1）不成立（2）成立}
\end{problem}

\section{解答题}

\begin{problem}
已知 \(f(x)=\sin\left(\omega x+\frac{\pi}{3}\right)\)，\(\omega>0\).

\begin{enumerate}
\item 设 \(\omega=1\)，求 \(y=f(x)\)，\(x\in[0,\pi]\) 的值域；
\item \(a>\pi\)（\(a\in\R\)），\(f(x)\) 的最小正周期为 \(\pi\)，若在 \(x\in[\pi,a]\) 上恰有 \(3\) 个零点，求 \(a\) 的取值范围.
\end{enumerate}
\end{problem}

\begin{problem}
如图，\(PA\)，\(PB\)，\(PC\) 为圆锥三条母线，\(AB=AC\).

\begin{enumerate}
\item 证明：\(PA\perp BC\)；
\item 若圆锥侧面积为 \(\sqrt3\pi\)，\(BC\) 为底面直径，\(BC=2\)，求二面角 \(B-PA-C\) 的大小.
\end{enumerate}

\FigureLayoutDeclare{img/2024/shanghai_spring/q18_fig1.png}{8.5cm}{0bp 9.002bp 0bp 3.001bp}
\bitmapfigure[max width=0.34\linewidth]{img/2024/shanghai_spring/q18_fig1.png}
\end{problem}

\begin{problem}
水果分为一级果和二级果，共 \(136\) 箱，其中一级果 \(102\) 箱，二级果 \(34\) 箱.

\begin{enumerate}
\item 随机挑选两箱水果，求恰好一级果和二级果各一箱的概率；
\item 进行分层抽样，共抽 \(8\) 箱水果，求一级果和二级果各几箱；
\item 抽取若干箱水果，其中一级果共 \(120\) 个，单果质量平均数为 \(303.45\) 克，方差为 \(603.46\)；二级果 \(48\) 个，单果质量平均数为 \(240.41\) 克，方差为 \(648.21\)；求 \(168\) 个水果的方差和平均数，并预估果园中单果的质量.
\end{enumerate}
\end{problem}

\begin{problem}
在平面直角坐标系 \(xOy\) 中，已知点 \(A\) 为椭圆 \(\Gamma:\frac{x^2}{6}+\frac{y^2}{2}=1\) 上一点，\(F_1\)，\(F_2\) 分别为椭圆的左，右焦点.

\begin{enumerate}
\item 若点 \(A\) 的横坐标为 \(2\)，求 \(\left|AF_1\right|\) 的长；
\item 设 \(\Gamma\) 的上，下顶点分别为 \(M_1\)，\(M_2\)，记 \(\triangle AF_1F_2\) 的面积为 \(S_1\)，\(\triangle AM_1M_2\) 的面积为 \(S_2\)，若 \(S_1\ge S_2\)，求 \(\left|OA\right|\) 的取值范围；
\item 若点 \(A\) 在 \(x\) 轴上方，设直线 \(AF_2\) 与 \(\Gamma\) 交于点 \(B\)，与 \(y\) 轴交于点 \(K\)，\(KF_1\) 延长线与 \(\Gamma\) 交于点 \(C\)，是否存在 \(x\) 轴上方的点 \(C\)，使得 \(\overrightarrow{F_1A}+\overrightarrow{F_1B}+\overrightarrow{F_1C} =\lambda\left(\overrightarrow{F_2A}+\overrightarrow{F_2B}+\overrightarrow{F_2C}\right)\)（\(\lambda\in\R\)）
成立？若存在，请求出点 \(C\) 的坐标；若不存在，请说明理由.
\end{enumerate}
\end{problem}

\begin{problem}
记 \(M(a)=\{t\mid t=f(x)-f(a),x\ge a\}\)，\(L(a)=\{t\mid t=f(x)-f(a),x\le a\}\).

\begin{enumerate}
\item 若 \(f(x)=x^2+1\)，求 \(M(1)\) 和 \(L(1)\)；
\item 若 \(f(x)=x^3-3x^2\)，求证：对于任意 \(a\in\R\)，都有 \(M(a)\subseteq[-4,+\infty)\)，且存在 \(a\)，使得 \(-4\in M(a)\).
\item 已知定义在 \(\R\) 上的 \(f(x)\) 有最小值，求证“\(f(x)\) 是偶函数”的充要条件是“对于任意正实数 \(c\)，均有 \(M(-c)=L(c)\)”.
\end{enumerate}
\end{problem}

